\section{Use cases}\label{sec:usecases}

The receivers show the pattern bound in code. This section states, independently of any
language, the situations in which a domain is at risk once agents contribute, so that the
pattern can be applied to domains other than task boards. Each use case names a situation, the
threat, the invariant at stake and the teaching move that answers it. Table~\ref{tab:usecases}
lists twelve in six families and Figure~\ref{fig:coverage} places each on the ladder, both in the
receivers as first published and as they stand after the work described here.

\begin{table}[t]
  \centering\small
  \caption{Twelve use cases in six families. The last two columns give the rung at which the
  receivers asserted the invariant at commit \texttt{8c4a74a} and now.}
  \label{tab:usecases}
  \begin{tabularx}{\linewidth}{@{}l X X c c@{}}
    \toprule
    & Situation & Invariant at stake & Then & Now \\
    \midrule
    \multicolumn{5}{@{}l}{\textit{Placement: sprawl spreads one rule across many places}} \\
    UC1 & A memoryless agent adds a feature & every part lands in the layer that owns it & L3 & L4 \\
    UC2 & A side effect is added to a use case & use cases publish facts; reactions subscribe & L2 & L4 \\
    UC3 & A business rule drifts to the edge & routes translate; they never decide & L0 & L4 \\
    \multicolumn{5}{@{}l}{\textit{Contamination: convenience pulls infrastructure into the core}} \\
    UC4 & A framework type enters the domain & the domain imports only the standard library & L3 & L4 \\
    UC5 & An outside capability leaks into a service & each capability has one owning layer & L3 & L4 \\
    \multicolumn{5}{@{}l}{\textit{Contract integrity: the promises the domain makes}} \\
    UC6 & A new failure mode is introduced & every error has one public translation & L5 & L5 \\
    UC7 & A test double lies & every double passes the real contract & L3 & L3 \\
    \multicolumn{5}{@{}l}{\textit{Harness erosion: the protections themselves wear down}} \\
    UC8 & The agent edits a rule to go green & rules change only with a person's decision & L2 & L2 \\
    UC9 & A dependency is added & new dependencies need a person & L2 & L2 \\
    \multicolumn{5}{@{}l}{\textit{Cognitive load: what people and agents can hold in mind}} \\
    UC10 & A reviewer faces a large AI diff & intent is reviewable apart from code & L1 & L3 \\
    UC11 & An agent starts cold & the declaration stays small & L3 & L3 \\
    \multicolumn{5}{@{}l}{\textit{Transfer: taking the approach to a new domain}} \\
    UC12 & A new domain is embedded & each invariant is placed on the ladder deliberately & L1 & L1 \\
    \bottomrule
  \end{tabularx}
\end{table}

\begin{figure}[t]
  \centering
  % Use cases on the ladder: hollow circle at 8c4a74a, filled circle now, arrow where they moved.
\begin{tikzpicture}[font=\footnotesize, x=0.95cm, y=0.52cm]
  \foreach \l in {0,...,5} {
    \draw[muted!35] (\l,0.4) -- (\l,-12.4);
    \node[above, font=\footnotesize\bfseries, text=jotblue] at (\l,0.4) {L\l};
  }
  \foreach \y/\name/\then/\now in {
      1/{UC1 feature by a memoryless agent}/3/4,
      2/{UC2 side effect in a use case}/2/4,
      3/{UC3 business rule at the edge}/0/4,
      4/{UC4 framework in the domain}/3/4,
      5/{UC5 capability leaks}/3/4,
      6/{UC6 new failure mode}/5/5,
      7/{UC7 lying test double}/3/3,
      8/{UC8 agent edits a rule}/2/2,
      9/{UC9 dependency added}/2/2,
      10/{UC10 reviewer overload}/1/3,
      11/{UC11 cold start}/3/3,
      12/{UC12 new domain}/1/1} {
    \node[anchor=east] at (-0.5,-\y) {\name};
    \ifnum\then=\now\relax
      \fill[structure] (\now,-\y) circle (4.2pt);
    \else
      \draw[-{Stealth[length=2mm]}, structure, line width=0.8pt] (\then+0.18,-\y) -- (\now-0.2,-\y);
      \draw[structure, line width=0.9pt, fill=white] (\then,-\y) circle (4.2pt);
      \fill[structure] (\now,-\y) circle (4.2pt);
    \fi
  }
\end{tikzpicture}

  \caption{Where each use case sits on the ladder in the receivers, at commit \texttt{8c4a74a}
  (hollow) and now (filled). UC6 is at L5 in Java and TypeScript and at L3 in Python.}
  \label{fig:coverage}
\end{figure}

\subsection{Placement}

Placement failures are the signature of sprawl. In \textbf{UC1} an agent with no prior session
implements a feature end to end; the risk is not that the feature fails but that each layer gets
a plausible and inconsistent piece of it. The answer is a worked slice to copy, followed by
architecture rules to pass. \textbf{UC2} and \textbf{UC3} are subtler because the mistake
preserves behaviour. A notification called inline in a service, instead of published as an
event, produces the same log line; a work-in-progress check copied into a route produces the
same refusal. No behavioural test can tell the difference, and at commit \texttt{8c4a74a}
nothing in any receiver did. They were closed by two \pattern{Executable Rule}s, that only the
events layer may log or notify and that routes make no decisions, each with a
\pattern{Teaching Failure} naming where the code should go.

\subsection{Contamination}

In \textbf{UC4} a framework annotation looks convenient on a domain object; in \textbf{UC5} a
service needs a setting or a quick query. Both are answered by capability confinement
(Section~\ref{sec:receivers}), and both were already checked before this work; they gained
explained failures.

\subsection{Contract integrity}

\textbf{UC6}, a new failure mode, is the one use case the receivers answer at L5 in two of three
languages: the new error does not compile until it has a status. \textbf{UC7} is the test double
that lies, answered by an \pattern{Honest Double}. It has a history: in the first seeding run the
contract created tasks titled First, Second and Third, whose alphabetical order equals their
creation order, so a double that sorted by title could not be caught. Changing the titles to
Write, Review and Publish fixed it. The episode is a small instance of kaizen, and a reminder that
a check is only as good as its fixture.

\subsection{Harness erosion}

\textbf{UC8} and \textbf{UC9} concern the ladder itself. An agent whose build is red has an easy
way to make it green: relax the rule. No agent in our pilot did so, but only the manifest forbade
it (L2). Raising UC8 to L3 is a matter of repository settings rather than tests, for example by
requiring a named person's review for any change to the rule files. Dependencies (UC9) are
similar, and can be checked by pinning the dependency list in a test.

\subsection{Cognitive load}

\textbf{UC10} is the reviewer's problem. When a change adds a feature, the parts a person most
needs to judge are the new rule and the examples that illustrate it, not the plumbing. Keeping
requirements in their own file, with scenarios that cite them and a check that keeps the two in
step, puts the intent of a change in a small, separate diff. We claim this lowers the reviewer's
load; we have not yet measured it. \textbf{UC11} is the agent's problem: every session starts
cold, and guidance that grows to compensate ends up unread. Fading keeps the manifest under 500
words even after two new kinds of intent were added.

\subsection{Transfer}

\textbf{UC12} asks whether the pattern survives a move to a domain that is not a task board. Two
illustrations, not implemented, suggest how. In a double-entry ledger, the invariant that debits
equal credits can be stated in EARS as an unwanted behaviour, shown in a posting slice, checked by
a property test, and prevented by a transaction type that can only be built from balanced legs.
In medication dosing, the invariant that a dose may not exceed a weight-based maximum without a
clinician's override can be checked by boundary scenarios and prevented, in part, by types that
distinguish milligrams from milligrams per kilogram. In both, the question the ladder asks is the
same: how high must this rule climb, given what it costs to get it wrong?
